\documentclass[../comorg-part2.tex]{subfiles}
\begin{document}
\section{Memory sub-system}
\subsection{Organisation of memory}
\subsubsection{Functional classification}
\begin{figure}[H]
    \centering
    \includegraphics[width=0.8\textwidth]{system-storage-memory.png}
    \caption*{Functional view of computer memory}
\end{figure}
\begin{itemize}
    \item System memory contains the code and data actively being used by the system during run-time. The processor can execute code directly from system memory. Because it handles the active subset of a running program, it requires high speeds.
    \item Storage memory This represents the entirety of all code and data stored in a computer system. The processor cannot execute code directly from storage memory; the data must first be transferred into the system memory before the CPU can process it.
\end{itemize}

\subsubsection{Data retention classification}
\begin{itemize}
    \item Volatile memory is temporary; all stored data is completely lost the moment electrical power is removed. Because of its speed, it is typically used for system memory. The primary examples are Static RAM (SRAM) and Dynamic RAM (DRAM).
    \item Non-volatile memory provides permanent storage, retaining its data even when the electric power is cut off. It is typically used for main storage to keep operating systems, applications, and files safe long-term. Examples include Flash memory and magnetic hard disks.
\end{itemize}

\subsubsection{Memory trade-off}
\begin{figure}[H]
    \centering
    \includegraphics[width=0.8\textwidth]{memory-triangle.png}
    \caption*{Memory pyramid}
\end{figure}
\begin{itemize}
    \item Top of the pyramid represents memory that is physically closest to the CPU, such as internal Registers and Level 1/2 Cache. This memory has the fastest access times. However, it is also the most expensive to manufacture, meaning it has the smallest storage capacity.
    \item Moving down the pyramid, it gets less costly and they have much larger capacities than cache, but their access times are slower (ranging from nanoseconds for RAM down to milliseconds for HDDs).
    \item The bottom of the pyramid consists of offline and removable storage. These devices offer massive, highly scalable storage capacities at the lowest cost. The trade-off is that they have the slowest access times.
\end{itemize}

\subsection{Volatile Memory}
\subsubsection{Static Random Access Memory (SRAM)}
SRAM is optimized for speed and low access latency, making it the standard choice for processor caches (like L1 and L2 cache).\\
\textbf{Characteristics}: It is very fast and has low power consumption in both active and standby states. However, it requires a lot of physical space per bit, meaning it has low storage density. It is "static" because it holds its data indefinitely as long as power is applied, without needing maintenance.\\


\subsubsection{Dynamic Random Access Memory (DRAM)}

\end{document}